% Teil II: der durchgerechnete Weg.

\subsection*{Gegeben und gesucht}
$c>0$, $S>0$, $a\ge0$; $(a_n),(b_n),(c_n)$ positive Folgen mit
$a_n\to a$, $b_n\to\infty$, $c_n\le S$. Zu beweisen:
(a) $\lim(c\,a_n)=c\,a$; (b) $\lim(c\,b_n)=\infty$; (c) $\lim c_n/b_n=0$.

\subsection*{a) $\lim(c\,a_n)=c\,a$}
Zu $\varepsilon>0$: wegen $a_n\to a$ gibt es $N$ mit
$|a_n-a|<\varepsilon/c$ (da $c>0$) für alle $n>N$. Dann
\begin{equation}
  |c\,a_n-c\,a|=|c|\,|a_n-a|=c\,|a_n-a|<c\cdot\frac\varepsilon c
  =\varepsilon .
\end{equation}
Also $c\,a_n\to c\,a$. \checkmark

\subsection*{b) $\lim(c\,b_n)=\infty$}
Zu $M>0$: wegen $b_n\to\infty$ gibt es $N$ mit $b_n>\tfrac M c$ (da
$c>0$) für alle $n>N$. Dann
\begin{equation}
  c\,b_n>c\cdot\frac Mc=M .
\end{equation}
Also $c\,b_n\to\infty$. \checkmark

\subsection*{c) $\lim\frac{c_n}{b_n}=0$}
Zu $\varepsilon>0$: wegen $b_n\to\infty$ gibt es $N$ mit
$b_n>\tfrac S\varepsilon$ für alle $n>N$. Mit der Beschränkung
$c_n\le S$ ($c_n>0$):
\begin{equation}
  \left|\frac{c_n}{b_n}\right|=\frac{c_n}{b_n}\le\frac{c_n}{S/\varepsilon}
  <\frac{S}{S/\varepsilon}=\varepsilon .
\end{equation}
Also $c_n/b_n\to0$. \checkmark

\begin{ergebnis}
\begin{equation*}
  \lim(c\,a_n)=c\,a,\qquad
  \lim(c\,b_n)=\infty,\qquad
  \lim\frac{c_n}{b_n}=0 .
\end{equation*}
\end{ergebnis}

\begin{kontrolle}
Beispiel: $a_n=2+\tfrac1n$ ($a=2$), $c=3$: $c\,a_n=6+\tfrac3n\to6=3\cdot2$;
$b_n=n\to\infty$: $3n\to\infty$; $c_n=\sin n$ (beschränkt $\le1$):
$\tfrac{c_n}{n}\to0$. Passt zu allen drei Aussagen.
\end{kontrolle}

\begin{merksatz}
Ketten von ``zu jedem $\varepsilon$ existiert $N$'' beweisen, indem man
das Ziel $N$ aus der jeweiligen Prämissen (Konvergenz, Divergenz,
Beschränktheit) ableitet; dadurch folgt alles direkt die Definitionen.
\end{merksatz}
